\documentclass[11pt]{article}

\usepackage{amsmath, amssymb}
\usepackage{physics}
\usepackage{siunitx}
\usepackage{graphicx}
\usepackage{geometry}
\usepackage{booktabs}
\usepackage{enumitem}
\geometry{margin=1in}

\title{Tutorial on DC-Link Voltage Balancing Modulation\\
for Cascaded H-Bridge Converters}
\author{}
\date{}

\begin{document}

\maketitle

\section{Introduction}

This document provides a structured tutorial explanation of the modulation-based DC-link voltage balancing method for Cascaded H-Bridge (CHB) converters.

\section{System Description and Notation}

\subsection{Converter Structure}

A single-phase grid-connected CHB converter consists of:

\begin{itemize}
\item $N$ : Number of cascaded H-bridge cells
\item $V_{dc,xH}$ : DC-link voltage of cell $x$, $x = 1,2,\dots,N$
\item $v_{xH}(t)$ : AC output voltage of cell $x$
\item $v_g(t)$ : Grid voltage
\item $V_g$ : RMS grid voltage
\item $\omega = 2\pi f$ : Grid angular frequency
\item $i_g(t)$ : Grid current
\item $I_g$ : RMS grid current amplitude
\item $L$ : Grid-side filter inductance
\end{itemize}

Total converter output voltage:

\begin{equation}
v_{ac}(t) = \sum_{x=1}^{N} v_{xH}(t)
\end{equation}

\subsection{Cell Power}

Instantaneous power of cell $x$:

\begin{equation}
p_x(t) = v_{xH}(t)\, i_g(t)
\end{equation}

Average power over one grid period $T$:

\begin{equation}
P_x = \frac{1}{T} \int_0^T v_{xH}(t) i_g(t)\, dt
\end{equation}

If:

\[
P_1 \neq P_2 \neq \dots \neq P_N
\]

then DC voltages evolve as:

\begin{equation}
C_x \frac{dV_{dc,xH}}{dt}
=
\frac{P_{in,x} - P_x}{V_{dc,xH}}
\end{equation}

\section{Conventional PI-Based Voltage Balancing}

\subsection{Control Structure}

Three control loops:

\begin{enumerate}
\item Total DC voltage control
\item Grid current control
\item Individual DC-link balancing loops
\end{enumerate}

For each cell:

\begin{equation}
e_x = V_{dc,H}^* - V_{dc,xH}
\end{equation}

Balancing PI output:

\begin{equation}
\Delta v_x
\end{equation}

Final reference:

\begin{equation}
v_{xH}^* = v^* + \Delta v_x
\end{equation}

\subsection{Modulation Index}

\begin{equation}
M_x = \frac{V_{xH,1}}{V_{dc,xH}}
\end{equation}

Linear region:

\[
-1 \le M_x \le 1
\]

\subsection{Maximum Power Capability}

\begin{equation}
P_{max,conv} =
P_{nom}
\frac{V_{dc,nH}}{V_g}
\sqrt{
1 -
\left(
\frac{\omega L I_g}{N V_{dc,nH}}
\right)
}
\end{equation}

where:

\begin{itemize}
\item $P_{nom}$ : Nominal converter power
\item $V_{dc,nH}$ : DC-link voltage of one cell
\item $V_g$ : Grid voltage amplitude
\item $\omega$ : Grid angular frequency
\item $L$ : Filter inductance
\item $I_g$ : Grid current amplitude
\item $N$ : Number of cells
\end{itemize}

\subsection{Bandwidth Constraint}

\[
t_s < t_c < t_v < t_{bal}
\]

where:

\begin{itemize}
\item $t_s$ : Switching period
\item $t_c$ : Current loop period
\item $t_v$ : Voltage loop period
\item $t_{bal}$ : Balancing loop period
\end{itemize}

\section{Proposed Modulation-Based Method}

\subsection{Operating Modes}

Let:

\begin{itemize}
\item $v^*(t)$ : Reference sinusoidal modulation signal
\item $v_c^*(t)$ : Clamped modulation signal
\item $v_{nc}^*(t)$ : Non-clamped modulation signal
\end{itemize}

Clamped mode:

\begin{equation}
v_c^* =
\begin{cases}
1 & v_g > 0 \\
-1 & v_g < 0
\end{cases}
\end{equation}

Non-clamped mode:

\begin{equation}
v_{nc}^* = N v^* - v_c^*
\end{equation}

Constraint:

\begin{equation}
v_c^* + v_{nc}^* = N v^*
\end{equation}

\section{Sorting Algorithm}

At each balancing period:

\begin{enumerate}
\item Measure $V_{dc,xH}$
\item Sort voltages
\item Compare to reference $V_{dc,H}^*$
\item Assign clamped or non-clamped mode
\end{enumerate}

If:

\[
v_g i_g < 0
\]

\begin{itemize}
\item $V_{dc,xH} < V_{dc,H}^*$ → clamped
\item $V_{dc,xH} > V_{dc,H}^*$ → non-clamped
\end{itemize}

If:

\[
v_g i_g > 0
\]

Logic reverses.

\section{Fourier Analysis}

Fourier coefficients:

\begin{equation}
a_1 =
\frac{1}{\pi}
\int f(\theta)\cos\theta\, d\theta
\end{equation}

\begin{equation}
b_1 =
\frac{1}{\pi}
\int f(\theta)\sin\theta\, d\theta
\end{equation}

Fundamental magnitude:

\begin{equation}
F = \sqrt{a_1^2 + b_1^2}
\end{equation}

\subsection{Clamped Waveform}

\begin{equation}
f_c(\theta) =
\begin{cases}
1 & [0,\phi] \\
M\cos\theta & [\phi,\pi] \\
-1 & [\pi,\pi+\phi]
\end{cases}
\end{equation}

where:

\begin{itemize}
\item $\phi$ : Clamping angle
\item $M$ : Nominal modulation index
\end{itemize}

\subsection{Fundamental Component}

\begin{equation}
F_c(\phi) =
\frac{1}{\pi}
\sqrt{
\left(0.5M\sin(2\phi) - M\phi - 2\cos\phi + \pi M + 2\right)^2
+
\left(\sin\phi(2 - M\sin\phi)\right)^2
}
\end{equation}

Non-clamped fundamental:

\begin{equation}
F_{nc}(\phi) = 2M - F_c(\phi)
\end{equation}

\section{Performance Improvements}

\begin{itemize}
\item Imbalance capability increased by approximately 27\%
\item Dynamic response approximately 13 times faster
\item Ripple amplitude balanced across cells
\end{itemize}

\section{Conclusion}

The proposed modulation-based voltage balancing method:

\begin{itemize}
\item Avoids modulation saturation
\item Decouples balancing from outer-loop bandwidth
\item Uses waveform shaping instead of amplitude modification
\item Improves both dynamic performance and power imbalance capability
\end{itemize}

\end{document}